%%%%%%%%%%%%%%%%%%%%%%%%%%%%%%%%%%%%%%%%%
% Cleese Assignment (For Students)
% LaTeX Template
% Version 2.0 (27/5/2018)
%
% This template originates from:
% http://www.LaTeXTemplates.com
%
% Author:
% Vel (vel@LaTeXTemplates.com)
%
% License:
% CC BY-NC-SA 3.0 (http://creativecommons.org/licenses/by-nc-sa/3.0/)
% 
%%%%%%%%%%%%%%%%%%%%%%%%%%%%%%%%%%%%%%%%%

%----------------------------------------------------------------------------------------
%	PACKAGES AND OTHER DOCUMENT CONFIGURATIONS
%----------------------------------------------------------------------------------------

\documentclass[11pt]{article}

%%%%%%%%%%%%%%%%%%%%%%%%%%%%%%%%%%%%%%%%%
% Cleese Assignment
% Structure Specification File
% Version 1.0 (27/5/2018)
%
% This template originates from:
% http://www.LaTeXTemplates.com
%
% Author:
% Vel (vel@LaTeXTemplates.com)
%
% License:
% CC BY-NC-SA 3.0 (http://creativecommons.org/licenses/by-nc-sa/3.0/)
% 
%%%%%%%%%%%%%%%%%%%%%%%%%%%%%%%%%%%%%%%%%

%----------------------------------------------------------------------------------------
%	PACKAGES AND OTHER DOCUMENT CONFIGURATIONS
%----------------------------------------------------------------------------------------

\usepackage{lastpage} % Required to determine the last page number for the footer

\usepackage{graphicx} % Required to insert images

\setlength\parindent{0pt} % Removes all indentation from paragraphs

\usepackage[most]{tcolorbox} % Required for boxes that split across pages

\usepackage{booktabs} % Required for better horizontal rules in tables

\usepackage{listings} % Required for insertion of code

\usepackage{etoolbox} % Required for if statements

\usepackage[T1]{fontenc} % Output font encoding for international characters

\usepackage[utf8]{inputenc} % Required for inputting international characters

%----------------------------------------------------------------------------------------
%	MARGINS
%----------------------------------------------------------------------------------------

\usepackage{geometry} % Required for adjusting page dimensions and margins

\geometry{
	paper=a4paper, % Change to letterpaper for US letter
	top=3cm, % Top margin
	bottom=3cm, % Bottom margin
	left=2.5cm, % Left margin
	right=2.5cm, % Right margin
	headheight=14pt, % Header height
	footskip=1.4cm, % Space from the bottom margin to the baseline of the footer
	headsep=1.2cm, % Space from the top margin to the baseline of the header
	%showframe, % Uncomment to show how the type block is set on the page
}

%----------------------------------------------------------------------------------------
%	FONT
%----------------------------------------------------------------------------------------

\usepackage[utf8]{inputenc} % Required for inputting international characters
\usepackage[T1]{fontenc} % Output font encoding for international characters

\usepackage[sfdefault,light]{roboto} % Use the Roboto font

%----------------------------------------------------------------------------------------
%	HEADERS AND FOOTERS
%----------------------------------------------------------------------------------------

\usepackage{fancyhdr} % Required for customising headers and footers

\pagestyle{fancy} % Enable custom headers and footers

\lhead{\small\assignmentClass{ : }{}\ \assignmentTitle} % Left header; output the instructor in brackets if one was set
\chead{} % Centre header
\rhead{\small\ifdef{\assignmentAuthorName}{\assignmentAuthorName}{\ifdef{\assignmentDueDate}{Due\ \assignmentDueDate}{}}} % Right header; output the author name if one was set, otherwise the due date if that was set

\lfoot{} % Left footer
\cfoot{\small Page\ \thepage\ of\ \pageref{LastPage}} % Centre footer
\rfoot{} % Right footer

\renewcommand\headrulewidth{0.5pt} % Thickness of the header rule

%----------------------------------------------------------------------------------------
%	MODIFY SECTION STYLES
%----------------------------------------------------------------------------------------

\usepackage{titlesec} % Required for modifying sections

%------------------------------------------------
% Section

\titleformat
{\section} % Section type being modified
[block] % Shape type, can be: hang, block, display, runin, leftmargin, rightmargin, drop, wrap, frame
{\Large\bfseries} % Format of the whole section
{\thesection} % Format of the section label
{6pt} % Space between the title and label
{} % Code before the label

\titlespacing{\section}{0pt}{0.5\baselineskip}{0.5\baselineskip} % Spacing around section titles, the order is: left, before and after

%------------------------------------------------
% Subsection

\titleformat
{\subsection} % Section type being modified
[block] % Shape type, can be: hang, block, display, runin, leftmargin, rightmargin, drop, wrap, frame
{\itshape} % Format of the whole section
{(\alph{subsection})} % Format of the section label
{4pt} % Space between the title and label
{} % Code before the label

\titlespacing{\subsection}{0pt}{0.5\baselineskip}{0.5\baselineskip} % Spacing around section titles, the order is: left, before and after

\renewcommand\thesubsection{(\alph{subsection})}

%----------------------------------------------------------------------------------------
%	CUSTOM QUESTION COMMANDS/ENVIRONMENTS
%----------------------------------------------------------------------------------------

% Environment to be used for each question in the assignment
\newenvironment{question}{
	\vspace{0.5\baselineskip} % Whitespace before the question
	\section{} % Blank section title (e.g. just Question 2)
	\lfoot{\small\itshape\assignmentQuestionName~\thesection~continued on next page\ldots} % Set the left footer to state the question continues on the next page, this is reset to nothing if it doesn't (below)
}{
	\lfoot{} % Reset the left footer to nothing if the current question does not continue on the next page
}

%------------------------------------------------

% Environment for subquestions, takes 1 argument - the name of the section
\newenvironment{subquestion}[1]{
	\subsection{#1}
}{
}

%------------------------------------------------

% Command to print a question sentence
\newcommand{\questiontext}[1]{
	\textbf{#1}
	\vspace{0.5\baselineskip} % Whitespace afterwards
}

%------------------------------------------------

% Command to print a box that breaks across pages with the question answer
\newcommand{\answer}[1]{
	\begin{tcolorbox}[breakable, enhanced]
		#1
	\end{tcolorbox}
}

%------------------------------------------------

% Command to print a box that breaks across pages with the space for a student to answer
\newcommand{\answerbox}[1]{
	\begin{tcolorbox}[breakable, enhanced]
		\vphantom{L}\vspace{\numexpr #1-1\relax\baselineskip} % \vphantom{L} to provide a typesetting strut with a height for the line, \numexpr to subtract user input by 1 to make it 0-based as this command is
	\end{tcolorbox}
}

%------------------------------------------------

% Command to print an assignment section title to split an assignment into major parts
\newcommand{\assignmentSection}[1]{
	{
		\centering % Centre the section title
		\vspace{2\baselineskip} % Whitespace before the entire section title
		
		\rule{0.9\textwidth}{0.5pt} % Horizontal rule
		
		\vspace{0.75\baselineskip} % Whitespace before the section title
		{\LARGE \MakeUppercase{#1}} % Section title, forced to be uppercase
		
		\rule{0.9\textwidth}{0.5pt} % Horizontal rule
		
		\vspace{\baselineskip} % Whitespace after the entire section title
		\setcounter{section}{0}
	}
}

%----------------------------------------------------------------------------------------
%	TITLE PAGE
%----------------------------------------------------------------------------------------


\author{\textbf{\assignmentAuthorName}} % Set the default title page author field
\date{} % Don't use the default title page date field

\title{
	\textsc{\LARGE Haute École d'Ingénierie et de Gestion du Canton de Vaud}\\[1.5cm] % logo Ecole
	\includegraphics[scale=.2]{images/heig.png}\\[1cm] % Include a department/university logo - this will require the graphicx package
	\thispagestyle{empty} % Suppress headers and footers
	\vspace{0.1\textheight} % Whitespace before the title
	\textbf{\assignmentTitle}\\[4pt]
	\ifdef{\assignmentDueDate}{{\assignmentDueDate}\\}{} % If a due date is supplied, output it
	\ifdef{\assignmentClassInstructor}{{\large \textit{\assignmentClassInstructor}}}{} % If an instructor is supplied, output it
	\vspace{0.32\textheight} % Whitespace before the author name
}
\setcounter{page}{0} % Include the file specifying the document structure and custom commands

%----------------------------------------------------------------------------------------
%	ASSIGNMENT INFORMATION
%----------------------------------------------------------------------------------------

% Required
\newcommand{\assignmentQuestionName}{Question} % The word to be used as a prefix to question numbers; example alternatives: Problem, Exercise
\newcommand{\assignmentClass}{CNM} % Course/class
\newcommand{\assignmentTitle}{Acceleration of an AI application} % Assignment title or name
\newcommand{\assignmentAuthorName}{Kevin Jordil \& Olivier D'Ancona} % Student name


% Optional (comment lines to remove)
\newcommand{\assignmentClassInstructor}{Professor: Marina Zapater } % Intructor name/time/description
\newcommand{\assignmentDueDate}{Thursday,\ 02\ February\, 2023} % Due date

%----------------------------------------------------------------------------------------

\begin{document}

%----------------------------------------------------------------------------------------
%	TITLE PAGE
%----------------------------------------------------------------------------------------

\maketitle % Print the title page

\thispagestyle{empty} % Suppress headers and footers on the title page

\newpage

\assignmentSection{Stage 1 – Choosing an application}

\section{Multi-Layer Perceptron}

\paragraph*{Multi-Layer Perceptron (MLP)}is a type of artificial neural network commonly used for supervised learning tasks such as classification.
It consists of one input layer, some hidden layer and one output layer. In each layer, each neuron is connected to all neurons in the next layer.
The architecture that we built have 3 layers:

\paragraph*{Architecture} our neural network implementation is as follows:

\begin{itemize}
	\item Input layer of size 784 (28x28 images)
	\item Hidden layer of size 1000 (number of hidden neurons)
	\item Output layer of size 10 (number of class of the dataset)
\end{itemize}

We used a c implementation who consists of a single file with all the functions needed to train and test the neural network.

\paragraph*{Dataset} We used the fashion dataset which is a commonly used dataset in machine learning and computer vision task.
He consists of a collection of Zalando shopping items such as shirts, pants and shoes. There are in total 10 categories.
he class are T-shirt/top, Trouser, Pullover, Dress, Coat, Sandal, Shirt, Sneaker, Bag and Ankle boot.
Each image is 28 pixels in height and 28 pixels in width, for a total of 784 pixels in total.
Each pixel has a single pixel-value associated with it, indicating the lightness or darkness of that pixel, with higher numbers meaning darker.
This pixel-value is an integer between 0 and 255. The training and test data sets have 785 columns.
The first column consists of the class labels (see above), and represents the article of clothing.
The rest of the columns contain the pixel-values of the associated image.

\section{Grayscale Conversion Program}

\paragraph*{Definition} We ran into troubles with the cuda neural network acceleration so we choosed an auxiliary program to be accelerated with cuda.
The program converts a color image to a grayscale image which is a good idea to implement with our multi layer perceptron.
Because our mlp necessitates a grayscale image as input we can convert rgb images to grayscale images with this program.
Furthermore this task is a good example to show the acceleration of a program with cuda.

\paragraph*{Principle} For each pixel of the image, we calculate 0.56 * red + 0.33 * green + 0.11 * blue.
The result is the grayscale value of the pixel. We then set the red, green and blue values of the pixel to the grayscale value.

\pagebreak

\assignmentSection{Stage 2 – Analysing application bottlenecks}

\section{Execution time}

\paragraph*{MLP baseline} The baseline execution time is showed on table \ref{tab:mlp_baseline} with a total time of 575.676547s.

\begin{table}[h]	\centering
	\begin{tabular}{|c|c|c|c|c|c|}
		\hline
		\textbf{time} & \textbf{accuracy} & \textbf{error} & \textbf{samples/sec} & \textbf{gflop/s} \\ \hline
		67.270 s      & 40.99\%           & 0.045          & 743.28               & 2.20             \\ \hline
		134.516 s     & 61.63\%           & 0.022          & 743.53               & 2.20             \\ \hline
		201.754 s     & 74.50\%           & 0.019          & 743.64               & 2.20             \\ \hline
		268.999 s     & 82.60\%           & 0.016          & 743.54               & 2.20             \\ \hline
		336.225 s     & 87.68\%           & 0.015          & 743.77               & 2.20             \\ \hline
		403.459 s     & 90.91\%           & 0.013          & 743.67               & 2.20             \\ \hline
		470.688 s     & 93.01\%           & 0.013          & 743.73               & 2.20             \\ \hline
		537.911 s     & 94.43\%           & 0.012          & 743.79               & 2.20             \\ \hline
	\end{tabular}
	\caption{MLP Baseline execution time}
	\label{tab:mlp_baseline}
\end{table}

\paragraph*{MLP accelerated with OpenMP} The accelerated execution time is showed on table \ref{tab:mlp_accelerated} with a total time of 393.993599s.

\begin{table}[h]	\centering
	\begin{tabular}{|c|c|c|c|c|c|}
		\hline
		\textbf{time} & \textbf{accuracy} & \textbf{error} & \textbf{samples/sec} & \textbf{gflop/s} \\ \hline
		46.092 s      & 40.99\%           & 0.045          & 1084.79              & 3.21             \\ \hline
		92.241 s      & 61.63\%           & 0.022          & 1083.43              & 3.20             \\ \hline
		138.231 s     & 74.50\%           & 0.019          & 1087.20              & 3.22             \\ \hline
		184.198 s     & 82.59\%           & 0.016          & 1087.76              & 3.22             \\ \hline
		230.153 s     & 87.68\%           & 0.015          & 1088.01              & 3.22             \\ \hline
		276.135 s     & 90.91\%           & 0.013          & 1087.40              & 3.22             \\ \hline
		322.124 s     & 93.01\%           & 0.013          & 1087.23              & 3.22             \\ \hline
		368.109 s     & 94.43\%           & 0.012          & 1087.32              & 3.22             \\ \hline
	\end{tabular}
	\caption{MLP OpenMP accelerated execution time}
	\label{tab:mlp_accelerated}
\end{table}

\paragraph*{MLP accelerated with OpenMP and CUDA} The accelerated execution time is showed on table \ref{tab:mlp_accelerated_cuda} with a total time of 326.487s.

\begin{table}[h]	\centering
	\begin{tabular}{|c|c|c|c|c|c|}
		\hline
		\textbf{time} & \textbf{accuracy} & \textbf{error} & \textbf{samples/sec} & \textbf{gflop/s} \\ \hline
		64.577 s      & 36.72\%           & 0.061          & 774.27               & 2.29             \\ \hline
		130.425 s     & 57.61\%           & 0.034          & 759.32               & 2.25             \\ \hline
		195.328 s     & 70.85\%           & 0.029          & 770.39               & 2.28             \\ \hline
		261.072 s     & 79.38\%           & 0.025          & 760.52               & 2.25             \\ \hline
		326.487 s     & 84.74\%           & 0.024          & 764.35               & 2.26             \\ \hline
		391.392 s     & 88.17\%           & 0.021          & 770.37               & 2.28             \\ \hline
		455.994 s     & 90.38\%           & 0.021          & 773.96               & 2.29             \\ \hline
		519.951 s     & 91.87\%           & 0.019          & 781.79               & 2.31             \\ \hline
		584.651 s     & 92.86\%           & 0.020          & 772.79               & 2.29             \\ \hline
		649.924 s     & 93.56\%           & 0.018          & 766.02               & 2.27             \\ \hline
		714.872 s     & 94.03\%           & 0.018          & 769.84               & 2.28             \\ \hline
		780.575 s     & 94.48\%           & 0.017          & 761.00               & 2.25             \\ \hline
		845.612 s     & 94.74\%           & 0.016          & 768.80               & 2.27             \\ \hline
	\end{tabular}
	\caption{MLP OpenMP and CUDA accelerated execution time}
	\label{tab:mlp_accelerated_cuda}
\end{table}

\paragraph*{Analysis} as we can see, the OpenMP implementation is the fastest then the baseline and finally the CUDA and OpenMP is the slowest.
The CUDA profiler results are showed on table \ref{tab:mlp_cuda_profiler}.
The CUDA profiler shows that the CUDA memcpy is the most time consuming operation, followed by the backprop\_kernel and the CUDA memset.
So our kernel function is too short to benefit from the GPU acceleration. The cost to move the data to memory are too high.

\begin{table}[h]	\centering
	\begin{tabular}{|l|c|c|c|c|c|c|c|}
		\hline
		\textbf{Type} & \textbf{Time perc.} & \textbf{Time} & \textbf{Calls} & \textbf{Avg} & \textbf{Name}         \\ \hline
		GPU :         & 47.10\%             & 15.4163s      & 410754         & 37.531us     & CUDA memcpy HtoH      \\ \hline
		              & 35.72\%             & 11.6917s      & 68459          & 170.78us     & backprop\_kernel      \\ \hline
		              & 17.18\%             & 5.62353s      & 136918         & 41.072us     & CUDA memset           \\ \hline
		API calls:    & 74.08\%             & 145.662s      & 410754         & 354.62us     & cudaMemcpy            \\ \hline
		              & 13.21\%             & 25.9788s      & 136918         & 189.74us     & cudaMemset            \\ \hline
		              & 9.01\%              & 17.7185s      & 68459          & 258.82us     & cudaDeviceSynchronize \\ \hline
		              & 3.49\%              & 6.86282s      & 68459          & 100.25us     & cudaLaunchKernel      \\ \hline
		              & 0.18\%              & 344.80ms      & 7              & 49.257ms     & cudaMallocManaged     \\ \hline
		              & 0.04\%              & 72.091ms      & 68459          & 1.0530us     & cudaGetLastError      \\ \hline
		              & 0.00\%              & 953.77us      & 7              & 136.25us     & cudaFree              \\ \hline
		              & 0.00\%              & 108.55us      & 97             & 1.1190us     & cuDeviceGetAttribute  \\ \hline
		              & 0.00\%              & 10.521us      & 1              & 10.521us     & cuDeviceTotalMem      \\ \hline
		              & 0.00\%              & 7.3430us      & 3              & 2.4470us     & cuDeviceGetCount      \\ \hline
		              & 0.00\%              & 3.8540us      & 2              & 1.9270us     & cuDeviceGet           \\ \hline
		              & 0.00\%              & 1.5100us      & 1              & 1.5100us     & cuDeviceGetName       \\ \hline
		              & 0.00\%              & 938ns         & 1              & 938ns        & cuDeviceGetUuid       \\ \hline
	\end{tabular}
	\caption{MLP CUDA accelerated GPU profiling results}
	\label{tab:mlp_cuda_profiler}
\end{table}


\section{Complexity Analysis}

\section{Accelerated part}

\section{Theoretical acceleration}



\assignmentSection{Stage 3 – Acceleration}

\assignmentSection{Stage 4 – Analysis of results}

\end{document}
